\begin{table*}[t]
\small
\setlength{\tabcolsep}{6pt}
\renewcommand{\arraystretch}{1.02}
\centering
\begin{tabularx}{\textwidth}{@{}p{3.9cm}p{2.4cm}X@{}}
\toprule
\textbf{Field} & \textbf{Type} & \textbf{Description} \\
\midrule
\texttt{short\_description} & String & A short summary of the image content. \\

\texttt{objects} & Array of Objects &
List of prominent objects. For each object, we record the following fields:
a detailed \texttt{description} of the object, the object's \texttt{location}, \texttt{relative\_size}, \texttt{shape\_and\_color}, \texttt{texture} (opt.), \texttt{appearance\_details} (opt.),
\texttt{relationship} with other objects, and the object's \texttt{orientation} (opt.).
When a human is described, the following fields are also included: the human's \texttt{pose} (opt.), \texttt{expression} (opt.), \texttt{clothing} (opt.), \texttt{action} (opt.), \texttt{gender} (opt.) and \texttt{skin\_tone\_and\_texture} (opt.).
If the image is a cluster of objects then also includes the
\texttt{number\_of\_objects} (opt.).\\

\texttt{background\_setting} & String &
Description of the environment/background. \\

\texttt{lighting} & Object &
Illumination details, including
\texttt{conditions} (e.g., daylight, studio),
\texttt{direction} (key-light orientation), and
\texttt{shadows} (opt.). \\

\texttt{aesthetics} & Object &
Global visual properties, including
\texttt{composition} (e.g., rule of third, symmetrical),
\texttt{color\_scheme} (dominant palette) and
\texttt{mood\_atmosphere} (emotional tone). \\

\shortstack[l]{\texttt{photographic\_}\\\texttt{characteristics}} & Object (optional) &
Camera/focus parameters:
\texttt{depth\_of\_field},
\texttt{focus},
\texttt{camera\_angle}, and
\texttt{lens\_focal\_length}. \\

\texttt{style\_medium} & String &
Artistic/rendering medium (e.g., photograph, oil painting, 3D render). \\

\texttt{text\_render} & Array of Objects &
A list of prominent text renders in the image. For each text render it includes the
\texttt{text},
\texttt{location},
\texttt{size},
\texttt{color},
\texttt{font} and
\texttt{appearance\_details} (opt.). \\

\texttt{context} & String & Additional context that helps understand the image better. \\

\texttt{artistic\_style} & String & High-level artistic/stylistic description (e.g., cinematic, minimalism). \\
\bottomrule
\end{tabularx}
\caption{\textbf{Structure of the JSON-based captions.} Each field captures a distinct aspect of the image, ensuring semantic completeness and consistency.}
\label{tab:JSON-structure}
\end{table*}
